\SourceLeaf{437}%
\[
\begin{aligned}
d\xi&=\xi'd\lambda\\
d\eta&=\eta'd\lambda\\
d\zeta&=\zeta'd\lambda.
\end{aligned}
\]
Der Analogie nach nennen wir nun
\[
\frac{d\lambda}{ds}
\]
das Krümmungsmass der krummen Linie auf der krummen Fläche, und dessen Reciprocum
\[
\frac{ds}{d\lambda}
\]
den Krümmungshalbmesser; bezeichnen wir letztern durch $\varrho$, so ist
\[
\begin{aligned}
\varrho d\xi&=\xi'ds\\
\varrho d\eta&=\eta'ds\\
\varrho d\zeta&=\zeta'ds.
\end{aligned}
\]

Ist also unsere Linie eine kürzeste Linie, so werden $\xi',\eta',\zeta'$ den Grössen $X,Y,Z$ proportional sein müssen; allein, da zugleich
\[
\xi'\xi'+\eta'\eta'+\zeta'\zeta'=XX+YY+ZZ=1,
\]
so wird
\[
\xi'=\pm X,\qquad\eta'=\pm Y,\qquad\zeta'=\pm Z,
\]
und da ferner
\[
\begin{aligned}
\xi'X+\eta'Y+\zeta'Z&=\cos\lambda'L\\
&=\pm(XX+YY+ZZ)\\
&=\pm1
\end{aligned}
\]
wird und wir den Punkt $\lambda'$ immer so wählen, dass
\[
\lambda'L<90^{0}
\]
ist, so wird für die kürzeste Linie
\[
\lambda'L=0
\]
sein oder $\lambda'$ mit $L$ zusammenfallen müssen. Es ist also
